Seasonal variation in cardiovascular events is well established, with recurrent winter excess in myocardial infarction, heart failure, blood pressure, and related physiological measures. By contrast, direct evidence for seasonality in imaging-defined cardiac structure and function is much less mature. This distinction matters for cardiac MRI studies: if a seasonal signal exists in MRI-derived phenotypes, it may reflect a genuinely under-studied source of biological variation rather than a settled fact.

This project therefore had three goals: first, to review the literature on seasonality in heart size, shape, and function across imaging and closely related cardiovascular modalities; second, to review statistical methods for detecting seasonality in irregularly timed biomedical observations; and third, to translate those findings into a practical analysis framework for a cardiac MRI dataset containing approximately 1000 scans acquired at random times across three years.
